\section{CS231N 课程概览}

第二部分由 Ehsan Adeli 讲授，介绍了 CS231N 2025 年春季学期的课程架构和核心内容。

\subsection{什么是计算机视觉}

计算机视觉的核心目标是让机器``看见''并\textbf{理解}图像和视频。最基础的任务是\textbf{图像分类}（Image Classification）：给模型一张图像，模型输出一个类别标签。

\begin{figure}[H]
\centering
\includegraphics[width=0.85\textwidth]{slides-images/slide-061.jpg}
\caption{线性分类器的基本思想：在特征空间中找到一个超平面（linear function）来分离不同类别（如猫与狗），但线性模型在处理复杂数据时能力有限\protect\footnotemark}
\end{figure}
\footnotetext{Slides 第2页（Part 2）。}

\begin{knowledgebox}{从线性分类到深度学习}
线性分类器是最简单的方法：在特征空间中用一个超平面分隔不同类别。但当数据不是线性可分时，线性分类器就力不从心了。神经网络通过\textbf{堆叠多层非线性变换}来建模复杂的决策边界，这就是``深度''学习的含义——深度来自于多层级的表示学习。
\end{knowledgebox}

\subsection{课程四大主题}

课程内容组织为四大主题模块：

\subsubsection{主题一：深度学习基础}

前几周的内容涵盖：
\begin{itemize}
    \item 图像分类与线性分类器
    \item 损失函数与优化
    \item 正则化（Regularization）——控制模型复杂度，防止过拟合
    \item 神经网络的前向传播与反向传播
    \item 训练技巧与调试方法
\end{itemize}

\subsubsection{主题二：感知与理解视觉世界}

\begin{figure}[H]
\centering
\includegraphics[width=0.85\textwidth]{slides-images/slide-062.jpg}
\caption{计算机视觉的核心任务层级：从图像分类到语义分割、目标检测和实例分割\protect\footnotemark}
\end{figure}
\footnotetext{Slides 第3页（Part 2）。}

涵盖的核心视觉任务：
\begin{itemize}
    \item \textbf{语义分割}（Semantic Segmentation）：为每个像素分配类别标签（如草地、猫、天空），但不区分同类别的不同个体
    \item \textbf{目标检测}（Object Detection）：定位并分类图像中的物体，输出边界框（bounding box）
    \item \textbf{实例分割}（Instance Segmentation）：最精细的任务，结合检测和分割，每个物体实例拥有独立的掩码
    \item \textbf{视频理解}：视频分类、动作识别
    \item \textbf{多模态理解}：结合视觉、听觉等多种模态
    \item \textbf{可视化与可解释性}：理解模型关注图像的哪些区域
\end{itemize}

涵盖的核心模型架构：
\begin{itemize}
    \item \textbf{卷积神经网络}（CNN）：从图像到卷积、池化、全连接的处理流水线
    \item \textbf{循环神经网络}（RNN）：处理序列数据
    \item \textbf{Transformer 与注意力机制}：当前最先进的架构
\end{itemize}

\subsubsection{主题三：生成式与交互式视觉智能}

\begin{itemize}
    \item \textbf{自监督学习}（Self-supervised Learning）：从无标注数据中学习表示，是训练大规模模型的关键技术
    \item \textbf{生成模型}：包括风格迁移、文生图（Text-to-Image）
    \item \textbf{扩散模型}（Diffusion Models）：通过逆转逐步加噪过程来生成图像（课程作业三将实现一个从文字生成表情符号的扩散模型）
    \item \textbf{视觉语言模型}（Vision-Language Models）：在共享表示空间中连接文本和图像
    \item \textbf{3D 视觉}：从图像重建三维表示，对机器人和 AR/VR 至关重要
    \item \textbf{具身智能}（Embodied AI）：视觉驱动的智能体在物理世界中感知、规划和执行
\end{itemize}

\subsubsection{主题四：以人为中心的 AI}

\begin{itemize}
    \item \textbf{大规模分布式训练}：数据并行（Data Parallelism）、模型并行（Model Parallelism）等策略
    \item \textbf{AI 伦理与社会影响}：偏见、公平性、对人类生活的影响
\end{itemize}

\subsection{学习目标}

Ehsan Adeli 总结了本课程的核心学习目标：
\begin{enumerate}
    \item 将计算机视觉应用形式化为具体的\textbf{任务}
    \item 开发和训练\textbf{视觉模型}——处理图像和视频数据的深度学习模型
    \item 理解该领域的\textbf{现状与发展方向}
\end{enumerate}

\subsection{学术认可}

\begin{figure}[H]
\centering
\includegraphics[width=0.85\textwidth]{slides-images/slide-066.jpg}
\caption{深度学习领域的最高荣誉：2018年图灵奖授予 Geoffrey Hinton、Yoshua Bengio 和 Yann LeCun；2024年诺贝尔物理学奖授予 Geoffrey Hinton 和 John Hopfield\protect\footnotemark}
\end{figure}
\footnotetext{Slides 第7页（Part 2）。}

深度学习领域已获得最高等级的学术认可：
\begin{itemize}
    \item \textbf{2018年图灵奖}：授予 Geoffrey Hinton、Yoshua Bengio 和 Yann LeCun，表彰他们在深度学习领域的概念性和工程性突破
    \item \textbf{2024年诺贝尔物理学奖}：授予 Geoffrey Hinton 和 John Hopfield，表彰他们对神经网络的基础性贡献
\end{itemize}

\subsection{本章小结}

CS231N 2025年春季课程覆盖四大主题：深度学习基础、视觉感知与理解、生成式与交互式视觉智能、以人为中心的 AI。课程从最基础的线性分类器出发，逐步深入到 CNN、Transformer、扩散模型、视觉语言模型等前沿架构，最终涵盖 3D 视觉、具身智能等最新研究方向。

%% ========== 拓展阅读 ==========

\subsection{拓展阅读}

\begin{itemize}
    \item CS231N 课程官网：\url{http://cs231n.stanford.edu/}
    \item Hubel \& Wiesel 原始论文：``Receptive fields of single neurones in the cat's striate cortex'' (1959)
    \item David Marr, \textit{Vision: A Computational Investigation} (1982)
    \item MIT Summer Vision Project 原文：\url{https://dspace.mit.edu/handle/1721.1/6125}
    \item Rumelhart, Hinton \& Williams, ``Learning representations by back-propagating errors'' (Nature, 1986)
    \item Krizhevsky, Sutskever \& Hinton, ``ImageNet Classification with Deep Convolutional Neural Networks'' (NeurIPS, 2012)
    \item Deng et al., ``ImageNet: A Large-Scale Hierarchical Image Database'' (CVPR, 2009)
    \item Fei-Fei Li, ``How we teach computers to understand pictures'' (TED Talk)
\end{itemize}

%% ========== 总结与延伸 ==========

